% !TeX spellcheck = ru_RU
\documentclass[a4paper, 14pt, unknownkeysallowed]{extreport}
\usepackage[utf8]{inputenc}
\usepackage[T2A]{fontenc}
\usepackage[english,russian]{babel}
\usepackage{cmap}
\usepackage{enumitem}

\usepackage{csvsimple}
\usepackage{hyphenat} 
\usepackage{amstext, amsmath,amsfonts,amssymb,amsthm,mathtools} 

\usepackage{setspace}
\onehalfspacing 

\usepackage{geometry}
\geometry{left=30mm}
\geometry{right=10mm}
\geometry{top=10mm}
\geometry{bottom=20mm}

\usepackage{graphicx}
\graphicspath{{img/}} 

\usepackage{threeparttable}
\usepackage{bigdelim}

\usepackage{setspace}
\onehalfspacing % Полуторный интервал

\frenchspacing
\usepackage{indentfirst} % Красная строка

\usepackage{cases}

\usepackage[unicode,pdftex]{hyperref} % Ссылки в pdf
\hypersetup{hidelinks}

\usepackage{xcolor}
\usepackage{listings}
% Для листинга кода:
\lstset{%
	language=Matlab,   					% выбор языка для подсветки	
	basicstyle=\small\sffamily,			% размер и начертание шрифта для подсветки кода
	numbers=left,						% где поставить нумерацию строк (слева\справа)
	%numberstyle=,					% размер шрифта для номеров строк
	stepnumber=1,						% размер шага между двумя номерами строк
	numbersep=5pt,						% как далеко отстоят номера строк от подсвечиваемого кода
	frame=single,						% рисовать рамку вокруг кода
	tabsize=4,							% размер табуляции по умолчанию равен 4 пробелам
	captionpos=t,						% позиция заголовка вверху [t] или внизу [b]
	breaklines=true,					
	breakatwhitespace=true,				% переносить строки только если есть пробел
	escapeinside={\#*}{*)},				% если нужно добавить комментарии в коде
	backgroundcolor=\color{white},
}
\usepackage{caption}
\captionsetup{labelsep = endash}
\captionsetup[figure]{name = {Рисунок}, justification = centerlast}
\captionsetup[lstlisting]{justification = justified}
\captionsetup[table]{justification = justified}

\usepackage{titlesec}
\titleformat{\chapter}{\LARGE\bfseries}{\thechapter}{20pt}{\LARGE\bfseries}
\titleformat{\section}{\Large\bfseries}{\thesection}{20pt}{\Large\bfseries}

\renewcommand\labelitemi{\ --\ }
\renewcommand\labelenumi{\theenumi)}

\newcommand{\img}[3] {
	\begin{center}
		\begin{figure}[h!]
			\includegraphics[scale = #1]{img/#2}
			\caption{#3}
			\label{img:#2}
		\end{figure}
	\end{center}
}

\newcommand{\lst}[4]{
\lstinputlisting[language=Matlab, firstline=#1, lastline=#2, label=lst:#3,caption=#4]{../src/code.m}
}

\usepackage{siunitx,array,booktabs}
\usepackage{amsthm}
\usepackage{tabularx}
\newtheorem{theorem}{Theorem}
\newtheorem{definition}{Опредление}


\begin{document}
% !TeX spellcheck = ru_RU
\begin{titlepage}
	
	\newgeometry{pdftex, left=2cm, right=2cm, top=2.5cm, bottom=2.5cm}
	\fontsize{12pt}{12pt}\selectfont
	\noindent \begin{minipage}{0.15\textwidth}
		\includegraphics[width=\linewidth]{img/b_logo.jpg}
	\end{minipage}
	\noindent\begin{minipage}{0.9\textwidth}\centering
		\textbf{Министерство науки и высшего образования Российской Федерации}\\
		\textbf{Федеральное государственное бюджетное образовательное учреждение высшего образования}\\
		\textbf{«Московский государственный технический университет имени Н.~Э.~Баумана}\\
		\textbf{(национальный исследовательский университет)»}\\
		\textbf{(МГТУ им. Н.~Э.~Баумана)}
	\end{minipage}

	\noindent\rule{18cm}{3pt}
	\newline\newline
	\noindent ФАКУЛЬТЕТ $\underline{\text{«Информатика и системы управления»~~~~~~~~~~~~~~~~~~~~~~~~~~~~~~~~~~~~~~~~~~~~~~~~~~~~~~~}}$ \newline\newline
	\noindent КАФЕДРА $\underline{\text{«Программное обеспечение ЭВМ и информационные технологии»~~~~~~~~~~~~~~~~~~~~~~~}}$\newline\newline\newline\newline\newline\newline\newline
	
	
	\begin{center}
		\Large\textbf{Отчет по лабораторной работе № 4}
		
		\Large\textbf{по курсу <<Моделирование>>}
		
		\textbf{\newline}		
	\end{center}
	
	\noindent\textbf{Тема} $\underline{\text{~~Моделирование простейшей СМО~~~~~~~~~~~~~~~~~~~~~~~~~~~~~~~~~~~~~~~~~~~~~~~~~~~~~~~~~~~~~~~~~~~~~~~~}}$\newline\newline
	\noindent\textbf{Студент} $\underline{\text{~~Виноградов А. О.~~~~~~~~~~~~~~~~~~~~~~~~~~~~~~~~~~~~~~~~~~~~~~~~~~~~~~~~~~~~~~~~~~~~~~~~~~~~~~~~~~~~~~}}$\newline\newline
	\noindent\textbf{Группа} $\underline{\text{~~ИУ7-76Б~~~~~~~~~~~~~~~~~~~~~~~~~~~~~~~~~~~~~~~~~~~~~~~~~~~~~~~~~~~~~~~~~~~~~~~~~~~~~~~~~~~~~~~~~~~~~~~~~~~~}}$\newline\newline
	\noindent\textbf{Оценка (баллы)} $\underline{\textbf{~~~~~~~~~~~~~~~~~~~~~~~~~~~~~~~~~~~~~~~~~~~~~~~~~~~~~~~~~~~~~~~~~~~~~~~~~~~~~~~~~~~~~~~~}}$\newline\newline
	\noindent\textbf{Преподаватель} $\underline{\text{~~Рудаков И. В.~~~~~~~~~~~~~~~~~~~~~~~~~~~~~~~~~~~~~~~~~~~~~~~~~~~~~~~}}$\newline
	
	\begin{center}
		\vfill
		Москва~---~\the\year
		~г.
	\end{center}
	\restoregeometry
	
\end{titlepage}
\setcounter{page}{2}

\chapter{Постановка задачи}
Цель работы: теоретическое изучение некоторых типовых законов распределения случайных величин.

Согласно 2 варианту по списку было выполнено изучение следующих законов распределения:
\begin{itemize}
	\item равномерное распределение на отрезке;
	\item нормальное распределение.
\end{itemize}


\chapter{Равномерное распределение}

Случайная величина имеет равномерное распределение на отрезке [a, b], если ее плотность распределения p(x) и функция распределения F(x) равны

\begin{equation}
	p(x) = \begin{cases}
		\frac{1}{b - a}, a\leq x\leq b;\\
		0, x < a или x > b.\\
	\end{cases},   
	F(x) =  \begin{cases}
		0, x < a;\\
		\frac{x-a}{b-a}, a\leq x\leq b;\\
		1, x > b.
	\end{cases}
\end{equation}

Графики плотности распределения p(x) и функции распределения F(x) для разных значений параметров a и b приведены на рисунках \ref{img:1}--\ref{img:2}.


\chapter{Нормальное распределение}

Случайная величина распределена по нормальному закону, или имеет нормальное распределение, если ее плотность

\begin{equation}
	\varphi_{m,\sigma}(x) = \frac{1}{\sigma \sqrt{2\pi}} e^{-\frac{(x-m)^2}{2\sigma^2}} (-\inf<m<+\inf, \sigma > 0).
\end{equation}

Функция нормального распределения имеет следующий вид:
\begin{equation}
	\Phi_{m, \sigma}(x) = \frac{1}{\sigma \sqrt{2\pi}} \int_{-\inf}^{x} e^{-\frac{(x-m)^2}{2\sigma^2}} dx.
\end{equation}

Графики плотности распределения $\varphi_{m,\sigma}(x)$ и функции распределения $\Phi_{m, \sigma}$ для разных значений параметров m и $\sigma$ представлены на рисунках \ref{img:21}--\ref{img:22}.
\end{document}